\chapter{Tests et Validation}

\section{Introduction}

Ce chapitre rend compte de la validation de VM Marketplace au regard des objectifs
définis dans l'introduction générale. Tous les tests ont été réalisés sur
l'infrastructure HCS MESRS de production (\textit{mesrscloud.rnu.tn}, région
\texttt{tn-global-1}). En l'absence de suite de tests automatisés, la validation
suit un protocole de tests manuels structuré par module fonctionnel, avec le
provisionnement cloud réel comme principale preuve de validation.

% ─────────────────────────────────────────────────────────────────────────────
\section{Environnement de Test}
% ─────────────────────────────────────────────────────────────────────────────

\begin{table}[H]
\centering
\caption{Environnement de test}
\label{tab:test_env}
\begin{tabularx}{\linewidth}{lX}
\toprule
\textbf{Composant} & \textbf{Configuration} \\
\midrule
Région HCS & \texttt{tn-global-1}, projet « Marketplace » \\
Réseau & VPC avec sous-réseau \texttt{47a8cd69-519b-45fe-845b-23704f01ace6} \\
Images OS testées & Ubuntu 22.04 LTS (UUID HCS) \\
Flavor maître & \texttt{c7n.xlarge.4} (4 vCPU, 16 Go RAM) \\
Flavor worker & \texttt{c7n.xlarge.4} (4 vCPU, 16 Go RAM) \\
Disque système & 50 Go SSD \\
Version K8s provisionnée & 1.29.15 \\
Navigateur de test & Chromium / Firefox \\
Hôte backend & localhost:8000 (Python 3.11, FastAPI) \\
Hôte frontend & localhost:3000 (Next.js 14) \\
Moteur IA & localhost:8001, Ollama llama3.1 (8B) \\
\bottomrule
\end{tabularx}
\end{table}

% ─────────────────────────────────────────────────────────────────────────────
\section{Tests d'Authentification}
% ─────────────────────────────────────────────────────────────────────────────

\begin{table}[H]
\centering
\caption{Résultats des tests d'authentification}
\label{tab:auth_tests}
\begin{tabularx}{\linewidth}{lXll}
\toprule
\textbf{TC} & \textbf{Cas de test} & \textbf{Attendu} & \textbf{Résultat} \\
\midrule
TC-A01 & Inscription avec e-mail et mot de passe valides & 201, utilisateur stocké en BD & \cmark \\
TC-A02 & Inscription avec e-mail déjà existant & Erreur 400 retournée & \cmark \\
TC-A03 & Connexion avec identifiants valides & 200, jeton JWT retourné & \cmark \\
TC-A04 & Connexion avec mauvais mot de passe & 401 Non autorisé & \cmark \\
TC-A05 & Accès à un endpoint protégé sans jeton & 401 Non autorisé & \cmark \\
TC-A06 & Accès à un endpoint protégé avec jeton expiré & 401 Non autorisé & \cmark \\
\bottomrule
\end{tabularx}
\end{table}

% ─────────────────────────────────────────────────────────────────────────────
\section{Tests de Provisionnement de VM}
% ─────────────────────────────────────────────────────────────────────────────

\begin{table}[H]
\centering
\caption{Résultats des tests de provisionnement de VM}
\label{tab:vm_tests}
\begin{tabularx}{\linewidth}{lXll}
\toprule
\textbf{TC} & \textbf{Cas de test} & \textbf{Attendu} & \textbf{Résultat} \\
\midrule
TC-V01 & Compléter l'assistant avec config valide + paiement & Session ID retournée & \cmark \\
TC-V02 & Le flux de progression SSE s'ouvre et envoie les étapes & 6 étapes diffusées & \cmark \\
TC-V03 & La VM apparaît dans la console HCS après apply & Instance ECS en état Running & \cmark \\
TC-V04 & IP publique de la VM accessible via SSH & Connexion root par clé SSH réussie & \cmark \\
TC-V05 & E-mail de confirmation reçu & E-mail avec détails VM livré & \cmark \\
TC-V06 & Tableau de bord Netdata accessible sur port 19999 & Interface Netdata chargée avec métriques & \cmark \\
TC-V07 & La VM apparaît dans le tableau de bord & Carte VM visible avec IP correcte & \cmark \\
TC-V08 & Supprimer une VM depuis le tableau de bord & \texttt{terraform destroy} exécuté ; VM retirée de la BD & \cmark \\
TC-V09 & Carte Stripe invalide rejetée & Erreur 400, aucun provisionnement lancé & \cmark \\
\bottomrule
\end{tabularx}
\end{table}

\screenshotplaceholder{Console HCS ECS affichant une instance VM provisionnée en état Running}

% ─────────────────────────────────────────────────────────────────────────────
\section{Tests de Recommandation IA}
% ─────────────────────────────────────────────────────────────────────────────

\begin{table}[H]
\centering
\caption{Résultats des tests de recommandation IA}
\label{tab:ai_tests}
\begin{tabularx}{\linewidth}{lXll}
\toprule
\textbf{TC} & \textbf{Cas de test} & \textbf{Attendu} & \textbf{Résultat} \\
\midrule
TC-IA01 & Décrire une charge de travail d'entraînement ML & Flavor haut niveau recommandé & \cmark \\
TC-IA02 & Décrire un serveur web statique & Flavor bas niveau (1--2 vCPU) & \cmark \\
TC-IA03 & Description vague (ex. : « quelque chose de rapide ») & Score de confiance faible ($<0,5$) & \cmark \\
TC-IA04 & Moteur IA hors ligne & Erreur 503 retournée par le backend & \cmark \\
TC-IA05 & Recommandation pré-remplit l'assistant & Assistant ouvert avec les specs correctes & \cmark \\
\bottomrule
\end{tabularx}
\end{table}

\screenshotplaceholder{Résultat de recommandation IA --- analyse de charge de travail et flavor recommandé pour un entraînement de machine learning}

% ─────────────────────────────────────────────────────────────────────────────
\section{Tests du Terminal SSH}
% ─────────────────────────────────────────────────────────────────────────────

\begin{table}[H]
\centering
\caption{Résultats des tests du terminal SSH}
\label{tab:ssh_tests}
\begin{tabularx}{\linewidth}{lXll}
\toprule
\textbf{TC} & \textbf{Cas de test} & \textbf{Attendu} & \textbf{Résultat} \\
\midrule
TC-S01 & Ouvrir le terminal SSH pour une VM en cours d'exécution & xterm.js connecté, invite shell affichée & \cmark \\
TC-S02 & Exécuter des commandes (ls, top, apt) dans le terminal & Commandes exécutées, sortie rendue & \cmark \\
TC-S03 & Terminal pour VM sans IP publique & Erreur d'affichée gracieusement & \cmark \\
TC-S04 & Fermer l'onglet du navigateur & WebSocket et canal SSH fermés côté serveur & \cmark \\
\bottomrule
\end{tabularx}
\end{table}

\screenshotplaceholder{Terminal SSH dans le navigateur affichant une session shell root interactive sur une VM HCS}

% ─────────────────────────────────────────────────────────────────────────────
\section{Tests du Cluster Kubernetes}
% ─────────────────────────────────────────────────────────────────────────────

\begin{table}[H]
\centering
\caption{Résultats des tests du cluster Kubernetes}
\label{tab:k8s_tests}
\begin{tabularx}{\linewidth}{lXll}
\toprule
\textbf{TC} & \textbf{Cas de test} & \textbf{Attendu} & \textbf{Résultat} \\
\midrule
TC-K01 & Compléter l'assistant cluster + paiement & Session ID retournée & \cmark \\
TC-K02 & Flux de progression SSE --- 6 étapes & Toutes les étapes atteignent l'état « complet » & \cmark \\
TC-K03 & Instance ECS maître créée sur HCS & ECS en état Running avec EIP & \cmark \\
TC-K04 & Instance ECS worker créée (sans EIP) & ECS en état Running, IP privée uniquement & \cmark \\
TC-K05 & Le worker peut joindre apt via le NAT maître & Paquets K8s installés sur le worker & \cmark \\
TC-K06 & \texttt{kubectl get nodes} affiche les deux nœuds Ready & STATUS : Ready pour tous les nœuds & \cmark \\
TC-K07 & \texttt{kubectl cluster-info} affiche l'URL du serveur API & Serveur API à l'EIP maître:6443 & \cmark \\
TC-K08 & Déployer nginx et passer à 3 réplicas & Tous les pods Running, service accessible & \cmark \\
TC-K09 & Télécharger le kubeconfig depuis le tableau de bord & Fichier YAML valide ; \texttt{kubectl} se connecte & \cmark \\
TC-K10 & Supprimer le cluster depuis le tableau de bord & terraform destroy terminé ; entrée BD supprimée & \cmark \\
\bottomrule
\end{tabularx}
\end{table}

\paragraph{Sortie de validation kubectl.}

La sortie suivante a été obtenue après téléchargement du kubeconfig et exécution de
\texttt{kubectl} sur le cluster provisionné :

\begin{lstlisting}[language=bash,
  caption={Sortie de kubectl get nodes sur le cluster provisionné},
  label={lst:kubectl_nodes}]
$ kubectl get nodes -o wide
NAME                    STATUS   ROLES          AGE   VERSION    INTERNAL-IP
k8s-cluster-master      Ready    control-plane  12m   v1.29.15   192.168.129.237
k8s-cluster-worker-1    Ready    <none>         9m    v1.29.15   192.168.129.112
\end{lstlisting}

\begin{lstlisting}[language=bash,
  caption={Sortie de kubectl cluster-info},
  label={lst:kubectl_info}]
$ kubectl cluster-info
Kubernetes control plane is running at https://193.95.30.249:6443
CoreDNS is running at https://193.95.30.249:6443/api/v1/namespaces/...
\end{lstlisting}

\screenshotplaceholder{kubectl get nodes --- les deux nœuds en statut Ready avec la version v1.29.15}

% ─────────────────────────────────────────────────────────────────────────────
\section{Évaluation des Objectifs}
% ─────────────────────────────────────────────────────────────────────────────

Le tableau~\ref{tab:objectives} évalue chacun des cinq objectifs initiaux au regard
des preuves de validation.

\begin{table}[H]
\centering
\caption{Évaluation des objectifs}
\label{tab:objectives}
\begin{tabularx}{\linewidth}{lXl}
\toprule
\textbf{Objectif} & \textbf{Preuve} & \textbf{Statut} \\
\midrule
Application web full-stack avec assistant VM, Stripe et Terraform sur HCS
  & TC-V01 à TC-V08 validés ; VMs réelles provisionnées sur HCS de production & \cmark Atteint \\
Moteur de recommandation IA (Llama~3.1, local)
  & TC-IA01 à TC-IA05 validés ; recommandations générées en moins de 30 s & \cmark Atteint \\
Terminal SSH navigateur et supervision Netdata
  & TC-S01 à TC-S04 validés ; Netdata installé et accessible & \cmark Atteint \\
Provisionnement cluster K8s avec réseau pour workers privés
  & TC-K01 à TC-K10 validés ; cluster 2 nœuds Ready sur HCS & \cmark Atteint \\
Validation bout en bout sur HCS de production
  & Tous les tests réalisés sur \textit{mesrscloud.rnu.tn} & \cmark Atteint \\
\bottomrule
\end{tabularx}
\end{table}

% ─────────────────────────────────────────────────────────────────────────────
\section{Limites Identifiées}
% ─────────────────────────────────────────────────────────────────────────────

Les limites suivantes ont été identifiées et documentées dans le cadre d'une
évaluation honnête :

\begin{enumerate}
  \item \textbf{Suivi de progression en mémoire.} Les dictionnaires
        \texttt{progress\_store} et \texttt{cluster\_progress} sont stockés en mémoire
        de processus. Un redémarrage du backend pendant le provisionnement (qui dure
        5 à 10 minutes) orphelinise le flux de progression ; la VM ou le cluster sera
        tout de même créé sur HCS mais le frontend n'affichera plus de mises à jour.

  \item \textbf{Absence de suite de tests automatisés.} Le projet ne dispose pas de
        tests unitaires, d'intégration ni de pipeline CI. Toute la validation a été
        effectuée manuellement sur l'environnement HCS de production.

  \item \textbf{SQLite pour la persistance.} SQLite est adapté à un déploiement de
        développement à instance unique, mais ne supporte pas les écritures simultanées
        de plusieurs workers backend. Une migration vers PostgreSQL serait nécessaire
        pour une mise à l'échelle en production.

  \item \textbf{URLs localhost codées en dur.} Les variables \texttt{NEXT\_PUBLIC\_API\_URL}
        et les URLs du moteur IA sont codées en dur vers localhost. Les déploiements
        multi-hôtes ou containerisés nécessitent une configuration par variable
        d'environnement.

  \item \textbf{Absence de contrôle d'accès par rôle.} Il n'existe pas de RBAC
        (contrôle d'accès basé sur les rôles) ni de gestion de quota par utilisateur
        au niveau de la plateforme. Tout utilisateur authentifié peut provisionner
        autant de ressources que le quota du projet HCS le permet.

  \item \textbf{Sécurité du kubeconfig.} Le kubeconfig stocké en base de données
        accorde un accès cluster-admin. Il est transmis via HTTPS mais stocké en clair
        dans la base de données SQLite.
\end{enumerate}

% ─────────────────────────────────────────────────────────────────────────────
\section{Conclusion}
% ─────────────────────────────────────────────────────────────────────────────

Les cinq objectifs initiaux ont tous été atteints et validés sur l'infrastructure HCS
MESRS de production. La plateforme provisionne avec succès aussi bien des VM
individuelles que des clusters Kubernetes multi-nœuds, offre une recommandation
assistée par IA, et propose un accès SSH intégré ainsi qu'une supervision automatisée.
Les limites identifiées sont bien comprises et constituent une feuille de route claire
pour une évolution vers une version de production.
